\documentclass[10pt,a4paper]{amsart}
\usepackage[margin=19mm]{geometry}
\usepackage{amsmath,amssymb,mathtools}
\usepackage[scr=rsfs,cal=euler]{mathalfa}
\usepackage{xfrac}
\usepackage[unicode,hidelinks,bookmarks=false]{hyperref}
\newcommand{\ie}{i.e.\ }
\newcommand{\cf}{cf.\ }
\newcommand{\resp}{respectively\ }
\newcommand{\Subsection}{\subsection}
\providecommand{\nobreakdash}{\nobreak\hbox{-}}
\setlength{\emergencystretch}{2em}
\multlinegap0pt
\usepackage{bm}
\usepackage{bbm}
\usepackage{dsfont}
\usepackage{xspace}
\usepackage{cancel}
\usepackage{mathtools}
\usepackage{amsthm}
\makeatletter
\@ifundefined{remark}{\newtheorem{remark}{Remark}}{}
\makeatother
\makeatletter
\newcommand{\E}{\operatorname*{\mathbf{E}}\ilimits@}
\renewcommand{\P}{\operatorname*{\mathbf{P}}\ilimits@}
\newcommand{\R}{\mathbf R}
\newcommand{\T}{\mathsf{T}}
\newcommand{\e}{\mathrm{e}}
\renewcommand{\ge}{\geqslant}
\renewcommand{\geq}{\geqslant}
\renewcommand{\le}{\leqslant}
\renewcommand{\leq}{\leqslant}
\expandafter\ifx\csname proof\endcsname\relax
\renewenvironment{proof}[1][\proofname]{\par\pushQED{\qed}\normalfont \topsep6\p@\@plus6\p@\relax\trivlist\item[\hskip\labelsep#1\@addpunct{}]\ignorespaces}{\popQED\endtrivlist\@endpefalse}
\fi
\expandafter\ifx\csname proofsketch\endcsname\relax
\newenvironment{proofsketch}
  {\begin{proof}[\proofsketchname]}
\fi
\newcommand{\thetastar}{\theta^\star}
\makeatother
\title[Beyond Modern Asymptotics for Log-Likelihood Ratios in Logis]{Beyond Modern Asymptotics for Log-Likelihood Ratios in Logistic Regression}
\author{Hugo Chardon, Reese Pathak, Nikita Zhivotovskiy}
\date{Source: arXiv:2608.02507v1 (CC BY 4.0)}
\begin{document}
\maketitle

\noindent\emph{Source and licence.} Beyond Modern Asymptotics for Log-Likelihood Ratios in Logistic Regression. Version arXiv:2608.02507v1 (CC BY 4.0), distributed under the Creative Commons Attribution 4.0 International licence, as recorded in the arXiv licence metadata for this exact version and shown on the abstract page https://arxiv.org/abs/2608.02507v1. Full rights notice: licenses/arxiv-2608.02507-CC-BY-4.0.txt.\par\smallskip

\noindent\textit{Editorial scope.} Counted unit: the Remark whose source label is \texttt{None} in the article (source0.tex lines 1106--1108, source label \texttt{None}), delivered together with the article's own complete original proof of it (source0.tex lines 1110--1205, one \texttt{proof} environment of 96 lines). No mathematics is written, repaired or re-derived: statement and proof are the article's own text line for line. Required context retained uncounted: none. Every definition, display and label of those lines is retained unchanged. Omitted from this package and not redistributed: 1--1105, 1109--1109, 1206--4708. Nothing inside a retained excerpt is omitted or altered: every retained body is delivered exactly as the article's lines are written, so no omission receipt of any kind accompanies this package. Citations: the retained excerpts carry no citation command at all, so no bibliography entry of the article is transcribed and no reference is redistributed. Reference adaptation: none was needed and none was made; every reference inside the delivered unit resolves inside it, so no unresolved reference or citation is produced. Printed slips kept literal: no printed word, symbol, hypothesis or number of the counted statement or of its proof was altered. Layout and class: the article's own class is \texttt{article} with packages including geometry, authblk, url, mathtools, xcolor, amsmath, amsthm, amsfonts, amssymb, framed, hyperref, enumitem, comment, xifthen; the wrapper is the lane's pinned \texttt{amsart} editorial wrapper at 10pt on A4, which restates the theorem-like environments the retained text needs and adds the article's own macro definitions that the retained text needs (\texttt{\textbackslash E}, \texttt{\textbackslash P}, \texttt{\textbackslash R}, \texttt{\textbackslash T}, \texttt{\textbackslash e}, \texttt{\textbackslash ge}, \texttt{\textbackslash geq}, \texttt{\textbackslash le}, \texttt{\textbackslash leq}, \texttt{\textbackslash thetastar}), with extra wrapper packages (bm, bbm, dsfont, xspace, cancel, mathtools). The delivered numbering is the unit's own. Assets: no retained span contains a figure environment, a graphics-inclusion command, a PDF-page inclusion, a table or a TikZ picture; no image file is copied into this package, the package declares no asset dependency and no image file is redistributed with it. Licence: version arXiv:2608.02507v1 of this article is distributed under the Creative Commons Attribution 4.0 International licence, as recorded in the arXiv licence metadata for this exact version and shown on the abstract page https://arxiv.org/abs/2608.02507v1. A full rights notice is delivered at licenses/arxiv-2608.02507-CC-BY-4.0.txt. Characters: every retained excerpt is pure ASCII, so the delivered engine, XeLaTeX with its default Latin Modern text font, reports no missing character.
\begin{remark}
The approach used in the proof is quite general and goes beyond the special case of the Vandermonde design. We will see that the key quantity is the maximal diagonal entry of the corresponding orthogonal projector. We note, however, that the computation is still quite specific to $\thetastar = 0$, and in general it will not scale well in high dimensions: it might lead to suboptimal dependence on $d$ in higher dimension, but still allows us to remove any dependence on $n$ in our case.
\end{remark}

\begin{proof}
Since $\thetastar = 0$, we have
$
\prod_{i = 1}^n \sigma(\langle x_i, \thetastar\rangle) = 2^{-n}.$
Moreover, $Y = (Y_1, \ldots, Y_n)$ consists of independent Rademacher random variables. Using that for all real $t$,
\begin{equation}
	\sigma(t) = \frac{1}{1+\e^{-t}} = \frac 1 2 \cdot \frac{\e^{t/2}}{\cosh(t/2)} \, ,
\end{equation}
hence $\log(2\sigma(t)) = t/2 - \log\cosh(t/2)$. Define the subspace $V = \{(a + bi + ci^2)_{i = 1}^n : a, b, c \in \R\}$ and, for every $i \in [n]$ and $\theta \in \R^{d}$, let $v_i = \langle \theta, x_i \rangle/2$. Since $\sigma(0)=1/2$, we deduce the following identity
\begin{align}
\Lambda_n^{\log}(0)
	= \sup_{\theta \in \R^{d}} \sum_{i=1}^n
		\log\Big(\frac{\sigma(Y_i\langle x_i,\theta\rangle)}{\sigma(0)} \Big)
	= \sup\limits_{v \in V}\left\{\langle Y, v\rangle - \sum\limits_{i = 1}^n \log(\cosh(v_i))\right\} \, .
\end{align}
 Let $\Pi_{V} \in \R^{n \times n}$ denote the orthogonal projector onto $V$. For every $v \in V$, we have
\[
\langle Y, v\rangle = \langle \Pi_{V} Y, v\rangle \leq \|\Pi_{V} Y\| \|v\|,
\]
and by the same computation
\begin{equation}
\label{eq:boundonvi}
|v_i| \leq |\langle \Pi_{V} e_i, v\rangle| \leq \|\Pi_{V} e_i\| \|v\| = \sqrt{e_i^{\T}\Pi_{V}^2 e_i}\,\|v\| = \sqrt{(\Pi_{V})_{ii}}\,\|v\|.
\end{equation}
We can easily check that an orthogonal basis of $V$ is given by $(u_0)_i = 1$, $(u_1)_i = i - \frac{n+1}{2}$, $(u_2)_i = ((u_1)_i)^2 - \frac{n^2 - 1}{12}$. An elementary computation gives
\begin{equation}
\label{eq:projdiagonal}
(\Pi_{V})_{ii}
= \sum_{k = 0}^2 \frac{(u_k)_i^2}{\|u_k\|_2^2}
=
\frac1n
+\frac{12\left(i-\frac{n+1}{2}\right)^2}{n(n^2-1)}
+\frac{180\left(\left(i-\frac{n+1}{2}\right)^2-\frac{n^2-1}{12}\right)^2}{n(n^2-1)(n^2-4)}
\leq \frac9n.
\end{equation}
Now consider
$
\psi(x) = \frac{\log(\cosh(x))}{x^2},
$
for $x > 0$,
and set $M_n^2 = \max_i (\Pi_{V})_{ii}$. An elementary analysis shows that $\psi$ is decreasing on $(0,+\infty)$. Thus, by \eqref{eq:boundonvi}, we have
\[
\sum\limits_{i = 1}^n \log(\cosh(v_i))
\ge
\sum\limits_{i = 1}^n v_i^2 \psi\left(M_n \|v\|\right)
=
\|v\|^2 \psi\left(M_n \|v\|\right)
=
\frac{\log(\cosh(M_n \|v\|))}{M_n^2},
\]
and therefore, combining the above inequalities, we reduce the upper bound to a one-dimensional optimization by taking the supremum with respect to the norm of $v$:
\[
\Lambda_n^{\operatorname{log}}(0)
\le
\sup_{r \geq 0}\left(r \|\Pi_{V} Y\| - \frac{\log(\cosh(M_n r))}{M_n^2}\right)
=
\frac{1}{M_n^2} \sup\limits_{r \geq 0}\left(r M_n \|\Pi_{V} Y\| - \log(\cosh(r))\right).
\]
It is straightforward to check that for any $w \in [0,1]$, it holds that
$
\sup\limits_{r \geq 0}(w r - \log(\cosh(r))) \leq w^2.
$
Therefore, whenever $M_n \|\Pi_{V} Y\| \leq \frac{1}{2}$, it holds that
\[
\Lambda_n^{\operatorname{log}}(0) \leq \frac{M_n^2 \|\Pi_{V} Y\|^2}{M_n^2} = \|\Pi_{V} Y\|^2.
\]
Note that
$
\|\Pi_{V} Y\|^2 = \sum_{k = 0}^2 \left\langle Y, \frac{u_k}{\|u_k\|} \right\rangle^2.
$
Since for any $k \in \{0,1,2\}$ and $\lambda \in \R$ it holds that
$
\E\exp\left(\lambda \left\langle Y, \frac{u_k}{\|u_k\|} \right\rangle\right) \leq \exp(\lambda^2/2),
$
we have for any $t > 0$,
\[
\P(\|\Pi_{V} Y\|^2 \geq 3 t^2)
\le
\P\left(\max_{k \in \{0,1,2\}} \left|\left\langle Y, \frac{u_k}{\|u_k\|} \right\rangle\right| \geq t \right)
\le
6\exp(-t^2/2),
\]
which implies that, on the intersection of the event $\left\{M_n \|\Pi_{V} Y\| \leq \frac{1}{2}\right\}$ and an event of probability at least $1-\delta$,
\[
\|\Pi_{V} Y\|^2 \leq 6\log(6/\delta).
\]
Note that when $n \geq 216\log(6/\delta)$, on the event $\|\Pi_{V} Y\|^2 \leq 6\log(6/\delta)$ we have, using \eqref{eq:projdiagonal},
\[
\|\Pi_{V} Y\|^2 \leq 6\log(6/\delta) \leq \frac{n}{36} \leq \frac{1}{4M_n^2},
\]
and therefore the event $\left\{M_n \|\Pi_{V} Y\| \leq \frac{1}{2}\right\}$ also holds. This proves the statement for $n \geq 216\log(6/\delta)$. Otherwise, if $n < 216\log(6/\delta)$, we have
\[
\Lambda_n^{\operatorname{log}}(0) \leq \log\left(\frac{1}{2^{-n}}\right) \leq 216\log(2)\log(6/\delta) \leq 216\log(6/\delta).
\]
The claim follows.
\end{proof}

\end{document}
